\subsection{诱导表示}

我们沿用前一日中关于 H 上的测度 \(\beta\) 与 G/H 上的测度 \(\nu\) 的记号和约定，以及函数 \(\kappa\colon\mathrm{G}\to\mathbf{R}_{+}^{*}\)。

存在一个连续函数 \(\eta\)，它定义在 \(\mathrm{G} \times \mathrm{G} / \mathrm{H}\) 上且取值于 \(\mathbf{R}_{+}^{*}\) 中，使得

\[
\eta (x, y \mathrm{H}) = \frac {\kappa (x y)}{\kappa (x)}
\]

对所有 \((x,y)\in\mathrm{G}\times\mathrm{G}\) 成立，且 \(\boldsymbol{\gamma}_{\mathrm{G}/\mathrm{H}}(x)\nu=(y\mapsto\eta(x^{-1},y))\cdot\nu\) 对 \(x\in G\) 成立（INT, VII, 第 56 页, § 2, 第 5 目, 定理 2, c)）。

设 \(p \in [1, +\infty[\)，并设 \(\pi\) 是 H 在复希尔伯特空间 E 中的一个酉表示。

引理 6. --- 设 \(f \in \mathcal{K}_{\overline{\pi}}(\mathrm{G})\)。对每个 \(g \in \mathrm{G}\)，函数

\[
\widetilde {f} \colon x \mapsto \eta (g ^ {- 1}, x \mathrm{H}) ^ {1 / p} f (g ^ {- 1} x)
\]

（它映 G 到 E 中）属于 \(\mathcal{K}_{\overline{\pi}}(\mathrm{G})\)，且满足 \(\mathrm{N}_p(\widetilde{f}) = \mathrm{N}_p(f)\)。

可以毫无困难地验证 \(\tilde{f} \in \mathcal{K}_{\overline{\pi}}(\mathrm{G})\)。由于

\[
\mathrm{N} _ {p} (\widetilde {f}) ^ {p} = \int_ {\mathrm{G/H}} ^ {*} \| \widetilde {f} \| ^ {p} d \nu = \int_ {\mathrm{G/H}} ^ {*} \gamma_ {\mathrm{G/H}} (g) (\| f \| ^ {p}) (y) \eta (g ^ {- 1}, y) d \nu (y)
\]

 以及 \((y\mapsto \eta (g^{-1},y))\cdot \nu = \pmb{\gamma}_{\mathrm{G / H}}(g)\nu\)，我们得到 \(\mathrm{N}_p(\widetilde{f}) = \mathrm{N}_p(f)\)。

 由这个引理可知，存在一个连续等距表示 \(\widetilde{\pi}\)，它从 G 映到 \(\mathrm{L}_{\overline{\pi}}^{p}(\mathrm{G})\) 中，使得对 \(f\in \mathcal{K}_{\pi}(\mathrm{G})\) 与 \(g\in \mathrm{G}\)，元素 \(\widetilde{\pi}(g)f\) 是上面定义的函数 \(\widetilde{f}\) 的类。若 \(p=2\)，则该表示是酉表示。

定义 2. --- 我们称如此定义的 G 在空间 \(L_{\overline{\pi}}^{2}(G)\) 中的酉表示为 G 的由 H 的表示 \(\pi\) 相对于 \(\kappa\) 所诱导的酉表示。把它记作 \(\operatorname{Ind}_{\mathrm{H}}^{\mathrm{G}}(\pi, \kappa)\)，或简作 \(\operatorname{Ind}_{\mathrm{H}}^{\mathrm{G}}(\pi)\)。

 注. --- 1) 设 \(\varrho\) 是 H 在复希尔伯特空间 F 中的一个酉表示，并设 \(u\colon\pi\to\varrho\) 是一个 H-态射。对每个函数 \(f\in\mathcal{K}_{\overline{\pi}}(\mathrm{G})\)，用 \(v(f)\) 表示 G 到 F 中的函数 \(g\mapsto u(f(g))\)；它属于 \(\mathcal{K}_{\overline{\varrho}}(\mathrm{G})\)，且满足 \(\mathrm{N}_{p}(v(f))\leqslant\|u\|\mathrm{N}_{p}(f)\)。由 \(\mathcal{K}_{\overline{\pi}}(\mathrm{G})\) 到 \(\mathcal{K}_{\overline{\varrho}}(\mathrm{G})\) 中的把 \(f\) 映为 \(v(f)\) 的线性映射因此延拓为 \(\mathrm{Ind}_{\mathrm{H}}^{\mathrm{G}}(\pi)\) 到 \(\mathrm{Ind}_{\mathrm{H}}^{\mathrm{G}}(\varrho)\) 中的一个连续 H-态射，记作 \(\mathrm{Ind}_{\mathrm{H}}^{\mathrm{G}}(u)\)。我们有 \(\mathrm{Ind}_{\mathrm{H}}^{\mathrm{G}}(1_{\pi})=1_{\mathrm{Ind}_{\mathrm{H}}^{\mathrm{G}}(\pi)}\)。设 \(\sigma\) 是 H 的一个酉表示，并设 \(v\colon\varrho\to\sigma\) 是一个 H-态射；我们有 \(\mathrm{Ind}_{\mathrm{H}}^{\mathrm{G}}(v\circ u)=\mathrm{Ind}_{\mathrm{H}}^{\mathrm{G}}(v)\circ\mathrm{Ind}_{\mathrm{H}}^{\mathrm{G}}(u)\)。

换言之，把 \(\pi\) 映为 \(\mathrm{Ind}_{\mathrm{H}}^{\mathrm{G}}(\pi)\)、把 \(u\) 映为 \(\mathrm{Ind}_{\mathrm{H}}^{\mathrm{G}}(u)\) 的构造是一个从 H 的酉表示的范畴到 G 的酉表示的范畴的函子（参见 CAT, I, § 2, 编纂中）。

\begin{enumerate}
\def\labelenumi{\arabic{enumi})}
\setcounter{enumi}{1}
\tightlist
\item
  设 \(\kappa^{\prime}\colon \mathrm{G}\to \mathbf{R}_{+}^{*}\)，使得
\end{enumerate}

\[
\frac {\kappa^ {\prime} (x h)}{\kappa^ {\prime} (x)} = \frac {\Delta_ {\mathrm{H}} (h)}{\Delta_ {\mathrm{G}} (h)} = \frac {\kappa (x h)}{\kappa (x)}
\]

 对所有 \((x,h)\in\mathrm{G}\times\mathrm{H}\) 成立。设 \(\nu^{\prime}\) 是 G/H 上的拟不变测度 \((\kappa^{\prime}\cdot\mu)/\beta\)。函数 \(\kappa^{\prime}\kappa^{-1}\) 通过取商定义了一个连续函数 \(\xi\colon G/H\to R_{+}^{*}\)，使得 \(\nu^{\prime}=\xi\cdot\nu\)。用 \(\alpha\) 表示 \(\mathcal{K}_{\pi}(G)\) 的由 \(f\mapsto(\kappa^{\prime}\kappa^{-1})^{1/p}f\) 定义的自同态；它满足

\[
\int_ {\mathrm{G} / \mathrm{H}} ^ {*} \| f \| ^ {p} d \nu^ {\prime} = \int_ {\mathrm{G} / \mathrm{H}} ^ {*} \| \alpha (f) \| ^ {p} d \nu
\]

并且它使我们可以把空间 \(\mathcal{L}_{\pi}^{p}(\mathrm{G},\nu^{\prime})\) 与 \(\mathcal{L}_{\pi}^{p}(\mathrm{G},\nu)\)、以及空间 \(\mathrm{L}_{\pi}^{p}(\mathrm{G},\nu^{\prime})\) 与 \(\mathrm{L}_{\pi}^{p}(\mathrm{G},\nu)\) 等同起来。此外，\(\alpha\) 诱导表示 \(\mathrm{Ind}_{\mathrm{H}}^{\mathrm{G}}(\pi,\kappa^{\prime})\) 与 \(\mathrm{Ind}_{\mathrm{H}}^{\mathrm{G}}(\pi,\kappa)\) 的一个等距同构。

